\section*{Chapitre \ref{ch:b2:unicellular-diversity} --- Diversité des organismes unicellulaires}

\begin{solution}{exo:b2:unicellular-diversity:1}
\omterm{def:b2:unicellular-diversity:unicellular}{Unicellulaire}~: une cellule assure toutes les fonctions et vit seule~;
\omterm{def:b2:unicellular-diversity:unicellular}{colonial}~: des cellules identiques restent attachées mais chacune
pourrait vivre seule~; \omterm{def:b2:unicellular-diversity:unicellular}{pluricellulaire}~: plusieurs types cellulaires
interdépendants, seule la lignée germinale se reproduit. \emph{E.~coli}
et la levure sont \omterm{def:b2:unicellular-diversity:unicellular}{unicellulaires} (un bourgeon se détache et vit seul)~;
\emph{Gonium} est \omterm{def:b2:unicellular-diversity:unicellular}{colonial} (16 cellules identiques, chacune capable de
fonder une colonie)~; \emph{Volvox} est \omterm{def:b2:unicellular-diversity:unicellular}{pluricellulaire} (cellules
somatiques mortelles, quelques cellules germinales).
\end{solution}

\begin{solution}{exo:b2:unicellular-diversity:2}
$S/V = 3/r$~: \qty{6}{\micro m^{-1}} pour le coque,
\qty{0.06}{\micro m^{-1}} pour le \omterm{def:b2:unicellular-diversity:protist}{protiste}, soit un facteur cent.
Le coque respire plus vite par unité de volume~: toute partie de son
cytoplasme se trouve à une fraction de micromètre de la surface par
laquelle entrent le dioxygène et la nourriture, alors que l'intérieur
du \omterm{def:b2:unicellular-diversity:protist}{protiste} est alimenté par cent fois moins de surface par unité de
volume.
\end{solution}

\begin{solution}{exo:b2:unicellular-diversity:3}
Pas de noyau (un \omterm{def:b2:unicellular-diversity:prokaryote}{nucléoïde}) contre une enveloppe nucléaire~; ni
mitochondries ni endomembranes, la chaîne respiratoire se trouvant dans
la membrane plasmique, contre des organites~; un flagelle qui est une
hélice rigide de flagelline mise en rotation par des protons, contre un
cil $9+2$ qui se courbe grâce à la dynéine et à l'ATP~; également la
taille (\qty{1}{\micro m} contre des dizaines) et une paroi de
peptidoglycane. Une bactérie ne peut pas phagocyter~: sa paroi lui
interdit d'englober une particule, ce que fait une amibe avec ses
pseudopodes.
\end{solution}

\begin{solution}{exo:b2:unicellular-diversity:4}
Cils~: locomotion et balayage de la nourriture vers la bouche~; sillon
oral~: canalise les particules vers le cytopharynx~; vacuole
digestive~: digère les bactéries englobées grâce aux enzymes
lysosomales~; \omterm{prop:b2:unicellular-diversity:osmoregulation}{vacuole contractile}~: expulse l'eau qui entre par
osmose~; macronucleus~: le noyau fonctionnel, dont les centaines de
copies de gènes sont transcrites pour la vie courante~; micronucleus~:
le noyau germinal diploïde, qui sert à la méiose et à la conjugaison.
\end{solution}

\begin{solution}{exo:b2:unicellular-diversity:5}
$t = x^2/2D$~: $(10^{-6})^2/10^{-9} = \qty{1}{ms}$ pour la bactérie~;
$(5\times 10^{-4})^2/10^{-9} = \qty{250}{s}$, environ quatre minutes,
pour l'amibe. L'amibe ne peut pas compter sur la diffusion pour
distribuer les métabolites~: il lui faut la cyclose et un transport
assuré par des moteurs le long du cytosquelette, dont la bactérie se
passe.
\end{solution}

\begin{solution}{exo:b2:unicellular-diversity:6}
$Re = \rho vL/\eta$~: \emph{Paramecium} $10^3\times 10^{-3}\times
2\times 10^{-4}/10^{-3} = 0.2$~; bactérie $10^3\times 2\times
10^{-5}\times 2\times 10^{-6}/10^{-3} = 4\times 10^{-5}$~; nageur
$10^3\times 1\times 2/10^{-3} = 2\times 10^6$. À $Re \gg 1$, l'inertie
porte le nageur après le mouvement~; à $Re \ll 1$, la traînée visqueuse
arrête la bactérie en moins d'un nanomètre, si bien qu'elle n'avance
que tant qu'elle bat.
\end{solution}

\begin{solution}{exo:b2:unicellular-diversity:7}
$v_s = 2r^2\Delta\rho g/9\eta = 2\times 10^{-10}\times 100\times
9.81/(9\times 10^{-3}) = \qty{2.2e-5}{m/s}$~; descendre de \qty{50}{m}
prend $50/2.2\times 10^{-5} = \qty{2.3e6}{s}$, soit 27 jours. Diviser le
rayon par deux divise $v_s$ par 4~: 108 jours. Réduire $\Delta\rho$ à
\qty{20}{kg/m^3} la divise par 5~: 135 jours.
\end{solution}

\begin{solution}{exo:b2:unicellular-diversity:8}
$V = \tfrac{4}{3}\pi\times 100\times 25\times 25 =
\qty{2.6e5}{\micro m^3}$. En \qty{1800}{s}~: $2.6\times 10^5/1800 =
\qty{145}{\micro m^3/s}$~; comme $\qty{1}{\micro m^3} =
\qty{1e-15}{L}$, cela fait \qty{1.5e-13}{L/s}.
\end{solution}

\begin{solution}{exo:b2:unicellular-diversity:9}
Un doublement sur \qty{1}{mm} correspond à une hausse d'environ
\qty{0.07}{\%} par micromètre (car $2^{1/1000} - 1 = 0.0007$), donc
\qty{0.07}{\%} d'un bout à l'autre de la cellule. Une course de
\qty{1}{s} couvre \qty{20}{\micro m}~: \qty{1.4}{\%}. Pour résoudre une
différence de \qty{0.07}{\%}, la cellule devrait compter quelque
$(1/0.0007)^2 = 2\times 10^6$ molécules à chaque extrémité, bien plus
qu'il ne s'en fixe en une seconde sur ses quelques milliers de
récepteurs~; \qty{1.4}{\%} en demande environ \num{5000}, ce qui est
faisable. La comparaison temporelle transforme une mesure spatiale
impossible en une mesure possible, en laissant la course faire
l'intégration.
\end{solution}

\begin{solution}{exo:b2:unicellular-diversity:10}
Surface $4\pi r^2$ avec $r = \qty{250}{\micro m}$~: \qty{7.9e5}{\micro
m^2}~; volume cytoplasmique $\approx$ surface $\times$ \qty{2}{\micro m}
$= \qty{1.6e6}{\micro m^3}$~; rapport \qty{0.5}{\micro m^{-1}}. Pour
\emph{E.~coli}, un cylindre~: $S/V = 2\pi r L/\pi r^2 L = 2/r =
\qty{4}{\micro m^{-1}}$. Le géant n'est que huit fois moins bien loti
qu'\emph{E.~coli}, et non 500 fois, car tout point de son cytoplasme
est à moins de \qty{2}{\micro m} de la surface. La vacuole stocke du
nitrate, son accepteur d'électrons, pour qu'il puisse respirer le
sulfure dans le sédiment anoxique entre les rares brassages qui
apportent du nitrate (\cref{ch:b2:microbial-metabolism}).
\end{solution}

\begin{solution}{exo:b2:unicellular-diversity:11}
Une cellule ne peut pas nager et se diviser à la fois. Une colonie de
16 cellules qui cesse de nager pendant les heures de la division coule,
d'après la \omterm{thm:b2:unicellular-diversity:stokes}{loi de Stokes}, d'une distance négligeable (son rayon est de
quelques dizaines de micromètres, sa vitesse de sédimentation de
quelques micromètres par seconde)~; la colonie entière peut se diviser
sans rien perdre. Une sphère de \num{10000} cellules et de rayon
\qty{500}{\micro m} coulerait cent fois plus vite ---
des millimètres par seconde, des mètres en une heure --- hors de la
lumière~: il est rentable de garder la plupart des cellules en train de
nager pendant que quelques-unes se divisent. Les cellules germinales
sont grandes parce qu'elles doivent porter les ressources nécessaires
pour construire, par divisions successives et sans se nourrir, une
sphère fille entière de milliers de cellules~; c'est donc le soma qui
les approvisionne.
\end{solution}

\begin{solution}{exo:b2:unicellular-diversity:12}
Un clade est un ancêtre et tous ses descendants~; un grade est un
niveau d'organisation atteint par plusieurs lignées. Les \omterm{def:b2:unicellular-diversity:protist}{protistes} ---
algues vertes (avec les plantes terrestres), diatomées (avec les algues
brunes), ciliés (avec les dinoflagellés et \emph{Plasmodium}), amibes
(proches des animaux et des champignons) --- n'ont en commun que la
cellule eucaryote vivant seule, et leur dernier ancêtre commun est celui
de tous les eucaryotes, nous compris~: c'est un grade. «~\omterm{def:b2:unicellular-diversity:prokaryote}{Procaryote}~»
réunit bactéries et archées par ce qui leur manque~; les archées sont
plus proches des eucaryotes que des bactéries, si bien que le terme
désigne là encore un grade. Les cyanobactéries descendent d'un ancêtre
unique qui a inventé la photosynthèse oxygénique et comprennent tous ses
descendants (la lignée des chloroplastes mise à part)~: c'est un clade.
\end{solution}

\begin{solution}{pb:b2:unicellular-diversity:1}
\textbf{1.} $1\times 1\times 0.1 = \qty{0.1}{mm^3} = \qty{0.1}{\micro L}$.
\textbf{2.} $24/\qty{0.1}{\micro L} = 240$ par \unit{\micro L} $=
\qty{2.4e5}{}$ cellules par millilitre.
\textbf{3.} $V = \tfrac{4}{3}\pi\times 125 = \qty{524}{\micro m^3}$~;
$S/V = 3/5 = \qty{0.6}{\micro m^{-1}}$.
\textbf{4.} $2.4\times 10^5\times 524 = \qty{1.26e8}{\micro m^3}$ par
millilitre, et $\qty{1}{mL} = \qty{1e12}{\micro m^3}$~: une fraction de
$1.3\times 10^{-4}$, soit \qty{0.013}{\%}.
\textbf{5.} $\qty{1000}{m^3} = \qty{1e9}{mL}$~: $2.4\times 10^{14}$
cellules.
\textbf{6.} La cyanobactérie est un filament de cellules bien plus
petites (quelques micromètres), bleu-vert, sans chloroplaste ni noyau
distincts, et elle ne nage pas à l'aide de flagelles~; l'algue est une
cellule ovoïde isolée, pourvue de deux flagelles, d'un stigma et d'un
chloroplaste unique en forme de coupe.
\textbf{7.} $384/24 = 16 = 2^4$~: quatre doublements en \qty{48}{h},
$T = \qty{12}{h}$.
\textbf{8.} $k = \ln 2/T = 0.693/12 = \qty{0.058}{h^{-1}}$.
\textbf{9.} $t = \ln(10^8/2.4\times 10^5)/k = \ln(417)/0.058 =
6.03/0.058 = \qty{104}{h}$, soit environ 4.3 jours.
\textbf{10.} Les nutriments (phosphate, nitrate) s'épuisent~; les
cellules se font de l'ombre, si bien que la lumière reçue par cellule
diminue~; les brouteurs (ciliés, rotifères, \emph{Daphnia}) se
multiplient à leurs dépens~; à \qty{1e8}{} cellules par millilitre, les
cellules occuperaient environ cinq pour cent du volume et épuiseraient le
$\mathrm{CO_2}$ dissous.
\textbf{11.} Croissance pendant la moitié du temps seulement~: le
rythme moyen est divisé par deux, \qty{208}{h}, soit 8.7 jours.
\textbf{12.} Le dénombrement mesure la population, non les cellules
individuelles~: si une cellule libère 8 cellules filles toutes les
\qty{36}{h}, la population est multipliée par 8 $= 2^3$ en trois \omterm{def:b2:unicellular-diversity:division}{temps
de génération} de \qty{12}{h} chacun. Le \omterm{def:b2:unicellular-diversity:division}{temps de génération} est le temps
moyen de doublement de l'effectif, quelle que soit la manière dont les
divisions sont groupées.
\textbf{13.} $Re = 10^3\times 10^{-4}\times 10^{-5}/10^{-3} = 10^{-3}$.
\textbf{14.} $m = 1050\times 5.24\times 10^{-16} = \qty{5.5e-13}{kg}$~;
$6\pi\eta r = 6\pi\times 10^{-3}\times 5\times 10^{-6} =
\qty{9.4e-8}{kg/s}$~; $d = 5.5\times 10^{-13}\times 10^{-4}/9.4\times
10^{-8} = \qty{5.8e-10}{m}$~: un demi-nanomètre.
\textbf{15.} $v_s = 2\times 25\times 10^{-12}\times 50\times
9.81/(9\times 10^{-3}) = \qty{2.7e-6}{m/s}$, soit \qty{2.7}{\micro m/s}.
\textbf{16.} Couler de \qty{1}{m}~: $1/2.7\times 10^{-6} =
\qty{3.7e5}{s}$, 4.3 jours~; remonter à la nage à \qty{100}{\micro m/s}~:
\qty{1e4}{s}, 2.8 heures. La nage l'emporte d'un facteur 37.
\textbf{17.} $c_\infty = \qty{1e-2}{mol/m^3}$~: $J = 4\pi\times
1.9\times 10^{-9}\times 5\times 10^{-6}\times 10^{-2} =
\qty{1.2e-15}{mol/s} = 7.2\times 10^8$ molécules par seconde.
\textbf{18.} \qty{100}{pg} $= \qty{1e-10}{g} = \qty{8.3e-12}{mol}$ de
carbone, doublés en \qty{12}{h} $= \qty{43200}{s}$~:
\qty{1.9e-16}{mol/s} $= 1.2\times 10^8$ atomes par seconde. Rapport de
l'apport à la demande~: $7.2/1.2 = 6$.
\textbf{19.} L'apport est multiplié par dix ($\propto r$), la demande
par mille ($\propto r^3$)~: le rapport devient $6/100 = 0.06$~; une
telle cellule ne pourrait croître qu'à six pour cent du rythme, si bien
qu'une algue de cette taille nourrie par diffusion est impossible sans
brassage, sans transport interne ou sans un métabolisme bien plus lent.
\textbf{20.} $\Delta\Pi = RT\Delta c = 8.314\times 293\times 99 =
\qty{2.4e5}{Pa}$, soit \qty{2.4}{bar}.
\textbf{21.} Surface $4\pi(5\times 10^{-6})^2 = \qty{3.14e-10}{m^2}$~;
flux entrant $L_p A\Delta\Pi = 10^{-14}\times 3.14\times 10^{-10}\times
2.4\times 10^5 = \qty{7.6e-19}{m^3/s} = \qty{0.76}{\micro m^3/s}$.
\textbf{22.} $524/0.76 = \qty{690}{s}$, environ onze minutes.
\textbf{23.} $P = \Delta\Pi\times$ débit $= 2.4\times 10^5\times
7.6\times 10^{-19} = \qty{1.8e-13}{W}$.
\textbf{24.} Un ATP~: $5\times 10^4/6.02\times 10^{23} =
\qty{8.3e-20}{J}$~: $1.8\times 10^{-13}/8.3\times 10^{-20} = 2.2\times
10^6$ ATP par seconde. Fixation du carbone~: $3\times 1.2\times 10^8 =
3.6\times 10^8$ ATP par seconde. La vacuole en coûte \qty{0.6}{\%}
--- une assurance bon marché contre l'éclatement.
\textbf{25.} Flux d'eau entrant \qty{0.76}{\micro m^3/s} (le volume de
la cellule en onze minutes)~; sans \omterm{prop:b2:unicellular-diversity:osmoregulation}{vacuole contractile}, la cellule
doublerait de volume en environ \qty{700}{s}~; maintenir son équilibre hydrique
coûte au moins $2\times 10^6$ ATP par seconde, moins d'un pour cent de
ce que sa photosynthèse dépense pour le carbone.
\end{solution}
